\documentclass[11pt]{article}
\usepackage[a4paper,margin=25mm]{geometry}
\usepackage{amsmath,amssymb,amsthm,hyperref}
\newtheorem{theorem}{Theorem}
\newtheorem{lemma}{Lemma}
\newtheorem{corollary}{Corollary}
\newcommand{\PP}{\mathcal P}
\title{Polynomially conditioned primitives in separated Legendre bands}
\author{Proof independently reviewed by root, size\_bounds, and independent\_directions}
\date{30 September 2026}
\begin{document}\maketitle
All unqualified norms use Lebesgue measure on $[0,1]$. Write
\[
 w(t)=t(1-t),\qquad \lambda_d=d(d+1),\qquad
 H_d(t)=\int_0^tP_d(2u-1)\,du.
\]
This note gives analytic tools for ordered-prefix moment constraints.
It does not assert independence of every family of such constraints,
or a construction of three simultaneous small dice.

\begin{lemma}[An exact primitive formula]\label{primitive}
Fix $d\ge2$ and $0\le r\le(d-2)/2$. The polynomial operator
\[
 \mathcal L_d C=(wC)''+\lambda_d C
\]
is invertible on $\PP_r$. For $Q\in\PP_r$, let
\[
 C_1=\mathcal L_d^{-1}Q',\qquad C_0=-\mathcal L_d^{-1}Q,
 \qquad B_a=w C_a.
\]
Then, identically on $[0,1]$,
\begin{align}
 \int_0^tQ(u)H_d'(u)\,du
   &=(Q-B_1')H_d+B_1H_d',\label{prim1}\\
 \int_0^tQ(u)H_d(u)\,du
   &=-B_0'H_d+B_0H_d'.\label{prim0}
\end{align}
The coefficient polynomials have degrees at most $r+2$ and norms
at most $C(r+1)^5\|Q\|_2$, with an absolute constant.
\end{lemma}
\begin{proof}
The shifted Legendre differential equation gives
$H_d''=-\lambda_d H_d/w$ in $(0,1)$.
For arbitrary $A,B=wC$,
\[
 (AH_d+BH_d')'=(A'-\lambda_d C)H_d+(A+B')H_d'.
\]
The two choices in the statement give the desired derivatives.
Both right sides vanish at zero because $H_d(0)=B_a(0)=0$.

For invertibility, use the shifted Jacobi polynomials with parameters
$(1,1)$, orthogonal for the weight $w$. The polynomial of degree $j$
satisfies
\[
 wJ_j''+2w'J_j'+j(j+3)J_j=0,
\]
so it is an eigenvector of $\mathcal L_d$ with eigenvalue
$\lambda_d-(j+1)(j+2)$. For $j\le r\le(d-2)/2$, these eigenvalues
are at least $c d^2$. Consequently
\[
 \|\mathcal L_d^{-1}F\|_{L^2(w)}\le C d^{-2}\|F\|_{L^2(w)}.
\]
For a polynomial $C$ of degree at most $r$,
$\|C\|_2\le C(r+1)\|C\|_{L^2(w)}$. Indeed, remove intervals
of length $c(r+1)^{-2}$ at the endpoints and use
$\|C\|_\infty\le(r+1)\|C\|_2$ to retain at least half its squared
norm; on the remainder $w\ge c(r+1)^{-2}$.
The polynomial $L^2$ Markov inequality
$\|Q'\|_2\le C(r+1)^2\|Q\|_2$ now bounds $C_1$ by
$C(r+1)^3d^{-2}\|Q\|_2$ and $C_0$ by
$C(r+1)d^{-2}\|Q\|_2$. Applying the same derivative estimate to
$wC_a$ proves the asserted, deliberately coarse coefficient bound.
\end{proof}

\begin{lemma}[Three separated bands]\label{three}
Suppose $d>2r+1$, $p,q,v\in\PP_r$, and $D\ge d+r+1$. Then
\[
 (\|p\|_2^2+\|q\|_2^2+\|v\|_2^2)^{1/2}
 \le C(D+1)^6\|p+H_dq+H_d'v\|_2.
\]
If $H_d,H_d'$ are both multiplied by $0<\eta\le1$, the bound
costs at most a further factor $\eta^{-1}$.
\end{lemma}
\begin{proof}
Put $G=H_dq+H_d'v$. Since $H_d$ is a linear combination of
shifted Legendre polynomials of degrees $d-1,d+1$, and $H_d'$ has
degree $d$ and is orthogonal to $\PP_{d-1}$, we have
$G\perp\PP_r$. For example, testing $H_dq$ against $p\in\PP_r$
reduces to orthogonality of $H_d$ against a polynomial of degree
at most $2r<d-1$. Thus $\|p+G\|_2^2=\|p\|_2^2+\|G\|_2^2$.

At the $d+1$ zeros $t_j$ of $H_d$, including the endpoints,
$G(t_j)=H_d'(t_j)v(t_j)$. The positive Gauss--Lobatto weights are
\[
 \omega_j=\frac1{d(d+1)H_d'(t_j)^2},\qquad \sum_j\omega_j=1,
\]
and this rule integrates polynomials through degree $2d-1$.
In particular $|H_d'(t_j)|\ge[d(d+1)]^{-1/2}$ and
\[
 \|v\|_2^2=\sum_j\omega_jv(t_j)^2
 \le d(d+1)\|G\|_\infty^2
 \le C(D+1)^4\|G\|_2^2.
\]
Here $\deg G\le D$ and the last step is the polynomial evaluation bound.
Since $\|H_d'\|_\infty\le1$,
$\|H_dq\|_2\le\|G\|_2+\|v\|_2$.
The reviewed multiplication estimate
$\|H_dq\|_2\ge c(D+1)^{-4}\|q\|_2$ from
\path{lacunary_legendre_conditioning.tex}
completes the proof. Scaling is immediate.
\end{proof}

\begin{theorem}[A fixed number of trinary bands]\label{many}
Fix $k$. Let $r\ge0$, $0<\eta_i\le1$, and positive integers $d_i$
satisfy
\[
 d_i>2\left(r+\sum_{j<i}(d_j+1)\right)+1.
\]
Set $D=r+\sum_{i=1}^k(d_i+1)$ and
$Z_{i,0}=1$, $Z_{i,1}=\eta_iH_{d_i}$, $Z_{i,2}=\eta_iH_{d_i}'$.
For arbitrary $p_\alpha\in\PP_r$, $\alpha\in\{0,1,2\}^k$,
\[
 \left(\sum_\alpha\|p_\alpha\|_2^2\right)^{1/2}
 \le C_k(D+1)^{6k}\prod_{i=1}^k\eta_i^{-1}
 \left\|\sum_\alpha p_\alpha\prod_{i=1}^kZ_{i,\alpha_i}\right\|_2.
\]
Minimal admissible choices have $D=O_k(r+1)$. Evaluation on any
node array whose star discrepancy is at most $c(D+1)^{-2}$ preserves
this lower bound up to an absolute constant.
\end{theorem}
\begin{proof}
Group the sum by its last coordinate. Each of its three coefficients
has degree at most $r+\sum_{j<k}(d_j+1)$, so Lemma~\ref{three}
applies. Iterate, summing the resulting squared inequalities. The
minimal degree recurrence grows by a factor at most three per step.
For sampling, every resulting polynomial $F$ has degree at most $D$,
and integration by parts for the empirical distribution function gives
\[
 \left|K^{-1}\sum_qF(b_q)^2-\int_0^1F^2\right|
 \le\operatorname{disc}(b)\int_0^1|(F^2)'|
 \le C(D+1)^2\operatorname{disc}(b)\|F\|_2^2.
\]
Choose the absolute discrepancy constant small enough.
\end{proof}

\paragraph{Use and limitation.}
An ordered prefix involving a new highest-frequency density contains
terms $\int QH_d'$ or $\int QH_d$, where the previous factors form a
lower-degree polynomial $Q$. Lemma~\ref{primitive} represents such a
term in two adjacent high-frequency bands with polynomially controlled
coefficients. Theorem~\ref{many} separates explicit trinary product
families. It does not automatically separate arbitrary ordered-prefix
families: after applying $\mathcal L_d^{-1}$, coefficient polynomials
need not retain a prescribed product form. Independence and positive
completion for those larger families remain separate questions.
\end{document}
